\section{System 0: The \texorpdfstring{$E_8$}{E8}-Persistence theory test case: Architectural requirements for zero-flux substrate}
\label{sec:System0_vacuum_requirements}
If the principles of IE are truly universal they must apply to the most fundamental layer of reality, the vacuum itself. In this section we will map the six components to this specific domain resulting in three non-negotiable requirements for the substrate of reality:
\begin{itemize}
    \item \textbf{Finiteness:} ($\mathrm{CAP}, \mathrm{MAR}$) To prevent energetic and informational divergence.
    \item \textbf{Unitarity:} ($\mathrm{MAP}, \mathrm{PRO}$) To ensure the lossless conservation of information.
    \item \textbf{Causality:} ($\mathrm{GOV}, \mathrm{TOL}$) To enforce a well-defined temporal evolution.
\end{itemize}

\subsection{The Closure Constraint}
\label{subsec:closure}
Within IE the vacuum is the minimal persistent substrate from which all observable structures emerge. Because no deeper persistent system exists beneath the vacuum, it cannot obtain configuration information, boundary conditions, synchronization signals, or memory from an external source and all information required to specify its state must be internally generated. We refer to this requirement as the \textbf{Closure Constraint}: the substrate of reality must be entirely self-defining. Its geometry must arise from its own architectural requirements. \textit{Any candidate substrate requiring outside data to specify its initial conditions is not a fundamental ground state.}

\subsection{Persistence Requires \texorpdfstring{$\mathrm{CAP}, \mathrm{MAR}$}{CAP, MAR} via Finiteness}
\label{sec:finiteness}

By the Closure Constraint the vacuum must be self-specifying. A self-specifying system cannot require infinite information to describe any local neighborhood as that would need an external decompressor that does not exist. Even with infinite storage, evaluating an infinitely complex model would take infinite time: each bit of the minimally compressed representation must be consulted to guarantee Governor containment, and any finite processing rate implies divergent latency. The system would be paralyzed before the next fluctuation arrives. Finiteness is therefore mandatory. The Capacity component enforces a finite state-space; the Margin component enforces finite reserves against fluctuation. Together we now show that finiteness require the substrate to be discrete, homogeneous, and a lattice.

\subsubsection{Finiteness Mandates a Discrete, Homogeneous Lattice}
\label{subsubsec:discrete}
\label{subsubsec:homogeneous}
\label{subsubsec:lattice}

A continuous volume admits infinite descriptive complexity. Under the Closure Constraint, no external algorithm exists to compress this to finite Kolmogorov complexity. Therefore, the substrate must be \textbf{Discrete}: there must exist a fundamental, indivisible unit of information that sets a maximum density. 

Furthermore, to minimize algorithmic processing overhead, the substrate must be \textbf{Homogeneous}. For the universe to be scalable with an $O(1)$ Kolmogorov Complexity independent of size, the local topology must remain constant regardless of location or direction. 

A structure that is both discrete (finite information density) and homogeneous (constant structure in every direction) is by definition a \textbf{Lattice}. This excludes continuous manifolds (infinite description) and random graphs (non-constant complexity). The lattice provides a discrete address space; it does not require the information payload residing at those addresses to be discrete. The Map and Protocol components may encode continuous state vectors within each node, subject only to the finiteness and unitarity constraints derived next.

\subsection{Persistence Requires \texorpdfstring{$\mathrm{MAP}, \mathrm{PRO}$}{MAP, PRO} via Unitarity}
\label{sec:unitarity}

For the vacuum to persist, it must maintain a stable Map over time. This requires that information is perfectly conserved. By the Closure Constraint (\cref{subsec:closure}), the vacuum must be a perfectly closed and lossless information network, making Unitarity a structural necessity, corresponding to the Map and Protocol components. This requirement imposes four structural properties on the substrate:

\subsubsection{Lattice Must Be Positive Definite}
\label{subsec:positive_definite}
\textit{Information-Theoretic Justification}: In a metric space, the norm $|\mathbf{v}|^2 = g_{\mu\nu}v^\mu v^\nu$ represents the information distance from the origin (reference state). For the system to have a well-defined minimum energy configuration (stable ground state), this distance must satisfy:
\begin{equation}
    |\mathbf{v}|^2 \geq 0 \quad \forall \mathbf{v} \neq 0
\end{equation}
A metric with negative eigenvalues (Lorentzian) allows $|\mathbf{v}|^2 < 0$, implying an imaginary ``distance'' from the reference state and preventing the definition of a stable ground state. Furthermore, a Lorentzian metric admits non-trivial null vectors ($|\mathbf{v}|^2 = 0$ where $\mathbf{v} \neq 0$), allowing for the creation of infinite information density without exceeding the capacity budget ($|k\mathbf{v}|^2 = 0$ for any scalar amplitude $k$), which violates the \textbf{Finiteness} component. Therefore, the \textbf{substrate} must be Euclidean.

\textbf{The Substrate-Projection Dual:}
The Euclidean positive-definiteness applies strictly to the substrate lattice. The observed Lorentzian signature $(-,+,+,+)$ emerges dynamically from the causal projection derived later in \cref{sec:metric_signature}; the substrate memory serves as a static, definite metric, while the active processor operates via an indefinite metric. This fundamental distinction provides the exact geometric mechanism by which time acquires its directional arrow, satisfying a core causal requirement of the persistence architecture.

\subsubsection{Lattice Must Be Mappable to a Torus}
\label{sec:torus}
To satisfy \textbf{Finiteness}, the system must be spatially bounded. To satisfy \textbf{Protocol} (Unitarity), it must be closed (no edges). By the combinatorial Gauss-Bonnet theorem, a regular lattice of fixed vertex coordination tiled on a closed surface of Euler characteristic $\chi$ requires topological defects if $\chi \neq 0$. The sphere ($\chi = 2$) therefore requires defects that break translational invariance. The Euclidean \textbf{Torus} ($T^n$, $\chi=0$) provides the unique flat, closed topology with periodic boundary conditions that admits a homogeneous, defect-free lattice structure, simultaneously satisfying Finiteness (bounded), Unitarity (closed), and Homogeneity (flat).

Consequently, the statistical evolution of the vacuum is defined by the Partition Function on a torus\footnote{Here, $\beta$ is the formal periodicity parameter of the imaginary time cycle. We use the partition function formalism for state enumeration on a torus, not to claim the vacuum has a literal temperature.}:
\begin{equation}
    Z(\beta) = \sum_{\text{states}} e^{-S_{\text{config}}} = \text{Tr}(e^{-\beta H})
\end{equation}

\subsubsection{Lattice Must Be Even}
\label{sec:even}
A stable Map requires path independence, which in turn requires evenness. A persistent system demands a unique, unambiguous equilibrium state. If the state count $Z$ depended on the path taken through moduli space (parameterization) rather than the configuration itself, the vacuum would lack a stable \textbf{Map}. 

This requires $Z(\beta)$ to be single-valued under the modular transformation $T: z \to z+1$. The Jacobi Theta Function, $\Theta_\Lambda(z) = \sum e^{i \pi z |\mathbf{v}|^2}$, counts these states. For \textbf{Even} lattices ($|\mathbf{v}|^2 = 2n$), the exponent $2\pi i n z$ is invariant under the shift. However, any odd-norm vector ($|\mathbf{v}|^2 = 2n+1$) introduces a sign inversion ($\Theta \to -\Theta$), rendering the Map multi-valued. Thus, the lattice must be \textbf{Even}.

\subsubsection{Lattice Must Be Self-Dual and Unimodular (Read/Write Symmetry)}
To prevent information loss via destructive interference or Landauer erasure, the encoding operation (write) and the decoding operation (read) must be informationally equivalent. This is enforced by the modular transformation $\mathcal{S}: z \to -1/z$, which maps the lattice to its reciprocal (Fourier dual).

By the Poisson Summation formula, the partition function transforms as:
\begin{equation}
    \Theta_\Lambda(-1/z) \propto \frac{1}{\text{vol}(\Lambda)} \Theta_{\Lambda^*}(z)
\end{equation}
For the Map to remain invariant ($Z_{\Lambda} = Z_{\Lambda^*}$), the lattice must be \textbf{Self-Dual} ($\Lambda = \Lambda^*$). This strictly enforces \textbf{Unimodularity} ($\text{vol}(\Lambda) = 1$), as the only scalar satisfying $V = 1/V$ is $1$. Any other volume introduces an irreversible scaling factor, violating Unitarity.

\textbf{Requirements:} The lattice of the vacuum must be \textbf{Positive Definite, Unimodular, Even, and Self-Dual}.

\subsection{Persistence Requires \texorpdfstring{$\mathrm{GOV}, \mathrm{TOL}$}{GOV, TOL} via Causality}
\label{sec:causality}
For the vacuum to persist, it must exist stably and \textit{evolve} in a well-defined manner. This creates a fundamental information-theoretic challenge: how to project the vast state space of a lattice node onto a single, linear temporal axis without ambiguity. 

\textbf{The Serialization Problem:} From the \textbf{Capacity} component, we establish that any persistent system must possess a finite state-space dimension, denoted $N$ (the number of linearly independent, orthogonal basis vectors spanning the node's total informational payload, i.e., the Hilbert-space dimension). To evolve, the system must serialize these $N$ potential states onto a discrete timeline defined by the \textbf{Temporal Modulus} ($\Delta$), representing the number of available temporal slots in one fundamental causal cycle.

\textbf{Causal Aliasing:} To express its full informational capacity, the system must serialize all $N$ potential states into $\Delta$ discrete temporal slots per fundamental cycle. By the Pigeonhole Principle, if $N > \Delta$, the number of states exceeds the available temporal slots. Consequently, at least one temporal coordinate must simultaneously host multiple distinct states. This results in \textbf{Causal Aliasing}: distinct information states become concurrently superimposed in time, destroying the system's ability to maintain a well-defined serial history.
    
\textbf{The Persistence Inequality:} To enforce a strict arrow of time and satisfy the Causality requirement, the projection must satisfy the Persistence Inequality: $\Delta \ge N$. The temporal container must meet or exceed the information content it holds. While the inequality allows $\Delta > N$, the Closure Constraint forbids inert auxiliary overhead. Therefore, any excess temporal capacity ($M = \Delta - N$) cannot be arbitrary; it must be explicitly mandated by additional structural requirements to maintain coherence.

\subsection{Summary: The Architectural Requirement of the Vacuum}
Through the application of Informational Energetics, we have fully specified the architectural requirements for the substrate of reality:

\begin{enumerate}
    \item ($\mathrm{CAP}, \mathrm{MAR}$) \textbf{Finite Lattice}, required by Finiteness.
    \item ($\mathrm{MAP}, \mathrm{PRO}$) \textbf{Positive Definite, Unimodular, Even, and Self-Dual}, required by Unitarity. 
    \item ($\mathrm{GOV}, \mathrm{TOL}$) Any causal system must satisfy a \textbf{Causal Projection}: a serialization of its $N$-dimensional state-space onto a discrete timeline of length $\Delta \ge N$.
\end{enumerate}

With this specification complete, the physical substrate is a constrained mathematical problem: find the unique geometric object that satisfies all of these requirements simultaneously.